\section{Diagrammi}\label{sec:diagrams}

In questa sezione discutiamo maggiormente la produzione dei diagrammi, oltre ai
risultati sulla \textit{connectedness}.\\
\\
Per produrre il diagramma UML (con profilo TinySOA) della Service Architecture a
partire dal BPMN, abbiamo seguito il metodo  SOMA (Service-Oriented Modeling and
Architecture). Prima di tutto, ogni pool rappresenta una \textit{capability},
modellata come un'\textit{Entity}. Ogni messaggio scambiato tra i partecipanti
indica un'operazione esposta dalla \textit{Service Interface} del
servizio/entità che riceve il messaggio --- per esempio, ACMEMobility permette
di avviare un noleggio.\\
Poiché il partecipante ACMEMobility costituisce il fulcro e l'orchestratore dei
processi, esso è stato rappresentato come un \textit{Task Service}, che può
usare o invocare gli altri (micro-) servizi, come dipendenze, con l'elemento
\textit{use} di UML --- interagendoci con protocolli quali \textit{http} e
\textit{SOAP}.\\
L'utente e i veicoli, essendo entità fisiche (umane o IoT), sono stati
rappresentati come semplici attori esterni che interagiscono con il sistema, in
contrapposizione alle \textit{Entity}, dei veri servizi, con delle
\textit{capability}, che agiscono ognuno su un dominio specifico. Inoltre, le
stazioni, essendo tante ma implementate allo stesso modo, sono rappresentate
come una \textit{box} invece che piatte.\\
Il risultante diagramma UML/TinySOA, dunque, combina le sintassi di
\textit{Class Diagram} (con servizi e interfacce) e \textit{Component Diagram}
(con più dettagli sulle entità) per mostrare, in un unico grafico,
l'architettura del sistema su più livelli contemporaneamente.\\



\subsection{Coreografia con sintassi formale}\label{sec:formal-choreography}

\textbf{Nota ---} Gli operatori di scelta ($+$) e di composizione parallela
($\mid$) hanno minore precedenza rispetto alla composizione sequenziale ($;$).\\

Si considerino i seguenti partecipanti:
\begin{itemize}
    \item \textbf{U}: User
    \item \textbf{A}: ACMEMobility
    \item \textbf{B}: Bank
    \item \textbf{F}: Fleet Management
    \item \textbf{S}: Station
    \item \textbf{V}: Vehicle
\end{itemize}

La coreografia $C$ si definisce come la scelta tra un noleggio immediato
($C_{imm}$) e una prenotazione ($C_{res}$):
\[ C = C_{imm} + C_{res} \]

\textbf{Noleggio immediato:}\\
\begin{align*}
C_{imm} = \\
& \ req\_imm: U \rightarrow A ; \\
& \ caut\_lock: A \rightarrow B ; ( \\
& \quad \quad caut\_lock\_err: B \rightarrow A ; C_{err\_reject} + \\
& \quad \quad caut\_locked: B \rightarrow A ; notify\_imm: A \rightarrow U ; C_{ride\_start} \\
& \ )
\end{align*}

\textbf{Prenotazione:}\\
\begin{align*}
C_{res} = \\
& \ req\_res: U \rightarrow A ; \\
& \ caut\_lock: A \rightarrow B ; ( \\
& \quad \quad caut\_lock\_err: B \rightarrow A ; C_{err\_reject} + \\
& \quad \quad caut\_locked: B \rightarrow A ; reserved: A \rightarrow U ; ( C_{cancel} + C_{scan\_res} ) \\
& \ )
\end{align*}

\textbf{Nota ---} Ogni richiesta alla banca può ritornare un \textit{ok} o un
\textit{errore}, nel qual caso si va a \textit{err\_manual}.

\begin{align*}
C_{cancel} = \\
& \ cancel: U \rightarrow A ; ( \\
& \quad \quad caut\_release: A \rightarrow B ; (caut\_release\_err: B \rightarrow A ; C_{err\_manual} + caut\_released: B \rightarrow A) + \\
& \quad \quad charge\_pen: A \rightarrow B; (charge\_pen\_err: B \rightarrow A ; C_{err\_manual} + charged\_pen: B \rightarrow A) \\
& \ ) \\
C_{scan\_res} = & \ scan\_res: U \rightarrow A ; C_{ride\_start}
\end{align*}

\textbf{Viaggio e inizio e fine del noleggio:}\\
\textbf{Nota ---} Anche il veicolo potrebbe ritornare successo o errore: per
semplicità, qui viene mostrato come se solo FleetManagement possa ritornare un
errore. Molti possibili errori come questo si potrebbero aggiungere in maniera
analoga ad altri presenti, ma complicherebbero le definizioni, aggiungendo tante
scelte e tanti livelli di annidamento.

\begin{align*}
C_{ride\_start} = \\
& \ track\_start: A \rightarrow F ; share\_start: F \rightarrow V ; share\_started: V \rightarrow F ; \\
& \ ( track\_err: F \rightarrow A; C_{err\_release} + \\
& \quad \quad track\_started: F \rightarrow A ; unlock: A \rightarrow S ; \\
& \quad \quad ( unlock\_err: S \rightarrow A ; C_{err\_release} + unlock\_ok: S \rightarrow A ; unlocked: A \rightarrow U ; C_{ride} ) \\
& \ )
\end{align*}

\textbf{Nota ---} Qui omettiamo per semplicità il loop in cui FleetManagement
riceve periodicamente dati (posizione e batteria) dal veicolo. Questo sarebbe un
ciclo separato che avviene in parallelo, dato che le informazioni vengono
memorizzate (in un database) per poter essere ritornate ad ACMEMobility su
richiesta, ma le sue formule sarebbero banalmente analoghe all'altro loop di
FleetManagement (ripetizione 0 o più volte di condivisione dati e risposta di
conferma).\\
\textbf{Nota ---} Si tenga presente che un tentativo senza successo di
parcheggio e blocco del veicolo può avvenire 1 o più volte, fino a un successo o
a una richiesta di assistenza: da coreografia, deve avvenire almeno una volta,
perché l'utente possa richiedere l'assistenza.\\

\begin{align*}
C_{ride} = \\
& \ (pos\_fetch: A \rightarrow F ; pos\_fetched: F \rightarrow A )^* ; \\
& \ (C_{park\_attempt})^* ; ( \\
& \quad \quad C_{park\_attempt} ; assist: U \rightarrow A + \\
& \quad \quad park: U \rightarrow A ; lock: A \rightarrow S ; lock\_ok: S \rightarrow A ; C_{ride\_end} \\
& \ ) \\
C_{park\_attempt} = & \ ( park: U \rightarrow A ; lock: A \rightarrow S ; lock\_err: S \rightarrow A ; park\_err: A \rightarrow U )
\end{align*}

\textbf{Nota ---} Ogni richiesta alla banca può ritornare \textit{ok} o
\textit{error}, nel qual caso si va a \textit{err\_manual}.\\
\textbf{Nota ---} Se fallisce la richiesta \textit{track\_stop}
(\textit{err\_track\_stop}), si va avanti, ma verrà richiesto un intervento
manuale in seguito.\\
\textbf{Nota ---} In questa parte conclusiva, nella maggior parte dei casi di
errore si usa \textit{notify\_err\_manual}, cioè si notifica l'utente, poiché,
nella coreografia, egli è in attesa o di un messaggio di fine corsa (con
successo, per esempio con una ricevuta) o di un messaggio di errore. Se l'errore
avviene al rilascio della cauzione, anche la coreografia non mostra notifiche
all'utente: si assume che il problema verrà risolto da un operatore al più
presto, senza però bloccare il processo.\\

\begin{align*}
C_{ride\_end} = \\
& \ ask\_batt: A \rightarrow F ; ( \\
& \quad \quad (err\_share\_batt: F \rightarrow A ; C_{err\_manual\_notify}) + \\
& \quad \quad share\_batt: F \rightarrow A ; \\
& \quad \quad track\_stop: A \rightarrow F ; share\_stop: F \rightarrow V , share\_stopped: V \rightarrow F ; \\
& \quad \quad ( track\_stopped: F \rightarrow A + (err\_track\_stopped: F \rightarrow A \mid C_{err\_manual})) ; \\
& \quad \quad charge: A \rightarrow B ; ( \\
& \quad \quad \quad \quad charge\_err: B \rightarrow A ; C_{err\_manual\_notify} + \\
& \quad \quad \quad \quad charged: B \rightarrow A ; receipt: A \rightarrow U ; caut\_release: A \rightarrow B ; ( \\
& \quad \quad \quad \quad \quad \quad caut\_release\_err: B \rightarrow A ; C_{err\_manual} + \\
& \quad \quad \quad \quad \quad \quad caut\_released: B \rightarrow A \\
& \quad \quad \quad \quad ) \\
& \quad \quad ))
\end{align*}

\textbf{Errori:}

Intervento manuale:\\

\begin{align*}
C_{err\_manual} = & \ 1 \\
\end{align*}

Intervento manuale, notificando l'utente:\\

\begin{align*}
C_{err\_manual\_notify} = & \ notify\_err\_manual: A \rightarrow U \\
\end{align*}

Non potendo iniziare la corsa, sblocca la cauzione e notifica l'utente, o
richiedi un intervento manuale anche qui, in caso di un errore della banca:\\

\begin{align*}
C_{err\_release} = \\
& \ caut\_release: A \rightarrow B ; ( \\
& \quad \quad C_{err\_manual} + caut\_released: B \rightarrow A ; C_{err\_reject} \\
& \ ) \\
\end{align*}

Rifiuta il noleggio subito, senza neanche coinvolgere la banca:

\begin{align*}
C_{err\_reject} = & \ notify\_reject: A \rightarrow U
\end{align*}


\subsection{Proiezioni}
Le proiezioni mostrano come la formula dell'intera coreografia si applica dal
punto di vista di un singolo partecipante.\\
Tutte le interazioni sono sincrone: la coreografia non ha ramificazioni
parallele. Nel BPMN, gli unici rami paralleli portano, come secondo percorso,
all'intervento manuale, che è molto semplice e non prosegue con ulteriori
comunicazioni con altri servizi.\\
\\
\textbf{Nota ---} Usiamo degli $1$ come placeholder, per indicare delle azioni
che non verrebbero codificate nella proiezione in questione, non comprendendo
una comunicazione tra i partecipanti --- per esempio, per gli errori manuali,
indicando che qualcosa succederà in quella sezione.\\
\textbf{Nota ---} Per alcuni partecipanti, l'intera coreografia o un suo
suffisso potrebbe non venire eseguito, in caso di errori occorsi nel
frattempo.\\
\textbf{Nota ---} Si ricorda che la sintassi \mbox{$proj(C, B) =
\overline{operation}@A$} indica la proiezione della coreografia $C$ sul
partecipante $B$, contenente l'operazione $operation$ con output verso il
partecipante $A$, poiché c'è una linea sopra; al contrario, $operation@A$ indica
un'operazione con input ricevuto da $A$.\\



\subsubsection{Proiezione per Bank (B)}
\begin{align*}
proj(C_{imm}, B) = & \ caut\_lock@A ; (\overline{caut\_lock\_err}@A + \overline{caut\_locked}@A ; proj(C_{ride\_start}, B) ) \\
proj(C_{res}, B) = \\
& \ caut\_lock@A ; ( \\
& \quad \quad \overline{caut\_lock\_err}@A + \\
& \quad \quad \overline{caut\_locked}@A ; (proj(C_{cancel}, B) + proj(C_{ride\_start}, B)) \\
& \ ) \\
proj(C_{cancel}, B) = \\
& \ caut\_release@A ; ( \\
& \quad \quad (\overline{caut\_release\_err}@A + \overline{caut\_released}@A) + \\
& \quad \quad charge\_pen@A ; (\overline{charge\_pen\_err}@A + \overline{charged\_pen}@A) \\
& \ ) \\
proj(C_{ride\_start}, B) = & \ caut\_release@A ; ( \overline{caut\_release\_err}@A + \overline{caut\_released}@A ; proj(C_{ride\_end}, B) ) \\
proj(C_{ride\_end}, B) = \\
& \ charge@A ; ( \\
& \quad \quad \overline{charge\_err}@A + \\
& \quad \quad ( \overline{charged}@A ; caut\_release@A ; ( \overline{caut\_release\_err}@A + \overline{caut\_released}@A ) ) \\
& \ )
\end{align*}

\subsubsection{Proiezione per Fleet Management (F)}
\begin{align*}
proj(C, F) = & \ ( \\
& \quad \quad 1 + \\
& \quad \quad track\_start@A ; share\_start@V ; \overline{share\_started}@V ; \\
& \quad \quad (\overline{track\_err}@A + \overline{track\_started}@A ; proj(C_{ride}, F) ) \\
& \ ) \\
proj(C_{ride}, F) = & \ (pos\_fetch@A ; \overline{pos\_fetched}@A )^* ; proj(C_{ride\_end}, F) \\
proj(C_{ride\_end}, F) = & \ ask\_batt@A ; ( \\
& \quad \quad \overline{err\_share\_batt}@A + \\
& \quad \quad \overline{share\_batt}@A ; track\_stop@A ; \\
& \quad \quad \overline{share\_stop}@V ; share\_stopped@V ; (\overline{err\_track\_stopped}@A + \overline{track\_stopped}@A ) \\
& \ )
\end{align*}

\subsubsection{Proiezione per ACMEMobility (A)}
\begin{align*}
proj(C_{imm}, A) = & \ req\_imm@U ; \overline{caut\_lock}@B ; \\
& \ ( caut\_lock\_err@B ; \overline{notify\_reject}@U + caut\_locked@B ; proj(C_{ride\_start}, A) ) \\
proj(C_{res}, A) = & \ req\_res@U ; \overline{caut\_lock}@B ; ( \\
& \quad \quad caut\_lock\_err@B ; \overline{notify\_reject}@U + \\
& \quad \quad caut\_locked@B ; \overline{notify\_reserved}@U ; (proj(C_{cancel}, A) + proj(C_{scan\_res}, A)) \\
& \ ) \\
proj(C_{cancel}, A) = & \ cancel@U ; ( \\
& \quad \quad \overline{caut\_release}@B ; (caut\_release\_err@B + caut\_released@B) + \\
& \quad \quad \overline{charge\_pen}@B ; (charge\_pen\_err@B + charged\_pen@B) \\
& \ ) \\
proj(C_{scan\_res}, A) = & \ scan\_res@U ; proj(C_{ride\_start}, A) \\
proj(C_{ride\_start}, A) = \\
& \ \overline{track\_start}@F ; ( \\
& \quad \quad track\_err@F ; proj(C_{err\_release}, A) + \\
& \quad \quad track\_started@F ; \overline{unlock}@S ; ( \\
& \quad \quad \quad \quad unlock\_err@S ; proj(C_{err\_release}, A) + \\
& \quad \quad \quad \quad unlock\_ok@S ; \overline{unlocked}@U ; proj(C_{ride}, A) \\
& \quad \quad ) \\
& \ ) \\
proj(C_{ride}, A) = \\
& \ {(\overline{pos\_fetch}@F ; pos\_fetched@F )}^* ; proj(C_{park\_attempt}^*, A) ; ( \\
& \quad \quad proj(C_{park\_attempt}, A) ; assist@U + \\
& \quad \quad park@U ; \overline{lock}@S ; lock\_ok@S ; proj(C_{ride\_end}, A) \\
& \ ) \\
proj(C_{park\_attempt}, A) = & \ park@U ; \overline{lock}@S ; lock\_err@S ; \overline{park\_err}@U \\
proj(C_{ride\_end}, A) = & \ \overline{ask\_batt}@F ; ( \\
& \quad \quad err\_share\_batt@F ; \overline{notify\_err\_manual}@U + \\
& \quad \quad \overline{share\_batt}@F ; \overline{track\_stop}@F ; ((err\_track\_stopped@F \mid 1) + track\_stopped@F) ; \\
& \quad \quad \overline{charge}@B ; ( \\
& \quad \quad \quad \quad charge\_err@B ; \overline{notify\_err\_manual}@U + \\
& \quad \quad \quad \quad charged@B ; \overline{receipt}@U ; \\
& \quad \quad \quad \quad \overline{caut\_release}@B ; ((caut\_release\_err@B ; 1) + caut\_released@B) \\
& \quad \quad ) \\
& \ ) \\
proj(C_{err\_release}, A) = & \ \overline{caut\_release}@B ; \\
& \ (caut\_release\_err@B ; 1 + caut\_released@B) ; \overline{notify\_reject}@U
\end{align*}

\subsubsection{Proiezione per Station (S)}

Si noti che la proiezione per la stazione include solo lo sblocco e il loop di
bloccaggio.\\

\begin{align*}
proj(C, S) = & \ ( \\
& \quad \quad 1 + \\
& \quad \quad unlock@A ; ( \\
& \quad \quad \quad \quad \overline{unlock\_err}@A + \\
& \quad \quad \quad \quad \overline{unlock\_ok}@A ; \\
& \quad \quad \quad \quad ( lock@A ; \overline{lock\_err}@A )^* ; ( lock@A ; \overline{lock\_ok}@A \mid 1 ) \\
& \quad \quad ) \\
& \ )
\end{align*}

\subsubsection{Proiezione per User (U)}
\begin{align*}
proj(C_{imm}, U) = & \ \overline{req\_imm}@A ; proj(C_{scan}, U) \\
proj(C_{res}, U) = & \ \overline{req\_res}@A ; ( \\
& \quad \quad notify\_reject@A + \\
& \quad \quad notify\_reserved@A ; (\overline{cancel}@A + \overline{scan}@A ; proj(C_{scan}, U)) \\
& \ )
\end{align*}

La proiezione \textit{scan} è valida sia per la scansione immediata che per
quella dopo una prenotazione:\\
\begin{align*}
proj(C_{scan}, U) = & \ notify\_reject@A + unlocked@A ; proj(C_{ride}, U) \\
proj(C_{ride}, U) = & \ (\overline{park}@A ; park\_err@A )^* ; \overline{park}@A ; ( \\
& \quad \quad park\_err@A ; \overline{assist}@A + \\
& \quad \quad park\_ok@A ; (notify\_err\_manual@A + receipt@A) \\
& \ )
\end{align*}

\subsubsection{Proiezione per Vehicle (V)}
\begin{align*}
    proj(C, V) = & \ ( \\
    & \quad \quad 1 + \\
    & \quad \quad share\_start@F ; \overline{share\_started}@F ; \\
    & \quad \quad ( 1 + share\_stop@F ; \overline{share\_stopped}@F ) \\
    & \ )
\end{align*}

Come per la definizione formale, sono stati omessi alcuni errori e il loop di
condivisione della posizione del veicolo.\\
Un esempio di loop di interazione tra FleetManagement e Vehicle sarebbe:\\
\begin{align*}
proj(C_{ride\_loop}, V) = & \ {( pos\_fetch@F ; \overline{pos\_fetched}@F )}^*
\end{align*}



\subsection{Analisi di di \textit{connectedness}}

Definiamo connessa una coreografia se le sequenze di interazioni garantiscono
una ``sicurezza causale'' (o ``causality safety''), ovvero se, in ogni momento,
un partecipante sa se sia il suo turno di agire, senza un controllo centrale che
glielo comunichi di volta in volta.\\
\\
Per l'analisi, possiamo considerare la nostra coreografia sincrona, dal momento
che ogni partecipante effettua una comunicazione in seguito ad un'altra, oppure
rimane ad aspettare il suo turno.\\
In realtà, la nostra architettura includerebbe due situazioni asincrone. La
prima è costituita dalle composizioni parallele; tuttavia, queste avvengono solo
all'estrema fine del processo e il secondo \textit{branch} è sempre molto
semplice, costituito dalla gestione manuale dell'errore ed eventualmente da una
comunicazione conclusiva di errore all'utente: dunque, la creazione di un
sotto-processo parallelo non rischia mai di interferire con il flusso delle
altre comunicazioni sincrone. La seconda eccezione è costituita dal loop che
avviene in background per la comunicazione delle informazioni di un veicolo
verso FleetManagement, che, come specificato più volte, abbiamo omesso dalle
modellazioni per semplicità, essendo banale. Difatti, esso comporta operazioni
completamente separate da quelle della coreografia, che non la intaccheranno in
alcun modo: FleetManagement continua solo ad accumulare dati nel suo database; a
una richiesta di ACMEMobility, ritornera gli ultimi disponibili.\\
\\
Considerando quindi la nostra coreografia sincrona, si può verificare la
\textit{connectedness} in quanto, per ogni sequenza \mbox{$o: a \rightarrow b ;
o': c \rightarrow d$}, è soddisfatta la condizione \mbox{$a=c \lor a=d \lor b=c
\lor b=d$}. Intuitivamente, ciò comporta che, in ogni coppia di comunicazioni
consecutive tra due servizi, almeno uno è coinvolto in entrambi: si può avere
una domanda e risposta tra gli stessi due servizi (\mbox{$a=d \land b=c$}), o
una serie di comunicazioni nella stessa direzione (\mbox{$a=c$}), oppure un
servizio che fa da ponte, ricevendo da uno e ritrasmettendo a un altro
(\mbox{$b=c$}).\\
In generale, quasi ogni passo del processo può essere ricondotto a una topologia
a stella in cui ACMEMobility agisce da orchestratore centrale: è presente in
quasi ogni comunicazione (come mittente o destinatario), facendo da ponte tra
l'utente e i vari microservizi che usa come dipendenze per ambiti specifici.
Esempi in cui non compare per qualche passaggio sono invece catene come quella
in $C_{ride\_start}$, in cui la ``parola'' passa, in ordine, per ACME,
FleetManagement, Vehicle, FleetManagement e nuovamente ad ACME:\@ ognuno
trasmette dopo aver ricevuto una comunicazione da qualcun altro.\\
\\
Per garantire la \textit{connectedness}, ci siamo accertati quindi
(eventualmente correggendo anche errori precedenti) che ogni servizio producesse
una risposta, come se dovesse passare il ``testimone'' del processo. È per
questo motivo che sia la coreografia che la definizione formale (e le
proiezioni) contengono molti messaggi di \textit{ok} o di \textit{errore},
altrimenti non si comprenderebbe come possa un servizio sapere l'esito della sua
richiesta e, dunque, come procedere. Alcuni messaggi di conferma mancano solo
alla fine del processo per un partecipante, per semplicità ma senza intaccare la
\textit{connectedness}: ad esempio, l'utente non aspetta una conferma dopo la
richiesta di cancellazione della prenotazione; se ACMEMobility avrà problemi con
l'addebito, provvederà in maniera asincrona con un'assistenza manuale.\\
\\
La \textit{connectedness}, in conclusione, fa sì che il processo possa sempre
andare avanti (\textit{liveness}) attraverso la semplice interazione tra i
partecipanti e senza un controllore centrale, fino a uno dei suoi stati di
terminazione.




%\cref{fig:host_performance_img} mostra spesso un picco in casa, ma suggerisce
%anche crescite iniziate prima.
%
%\begin{figure}[H] \centering
%    \includegraphics[width=\linewidth]{../../docs/img/Host_Performance_1996_2024.png}
%    \caption{Performance 1988--2024 per alcuni Pesi host}
%    \label{fig:host_performance_img} \end{figure}
